\section{第 19 轮：精确对偶、对偶证书的不可能性，与 Bloch 截面分解}\label{sec:r19}

\subsection{本轮要交代的三件事}\label{sec:r19:motiv}

第~\ref{sec:r18}~节把主线改写成 \eqref{eq:starmain}，并留下三个未了项：
命题~\ref{thm:concave}~的凹性当时只有数值证据；
推论~\ref{cor:dualcert}~因此是条件性的；
第八版队列的队首 P1（构造各向异性对偶证书 $S$）还没有动手。
本轮的处理是：凹性给出\textbf{解析证明}（定理~\ref{thm:concave}），
并顺带得到闭式梯度（推论~\ref{cor:gradpsi}）；
在此基础上证明 $\Lam$-参数化的对偶是\textbf{精确紧}的（定理~\ref{thm:exactdual}），
于是 P1 不是一个「还没做成」的任务，而是一个\textbf{不可能做成}的任务
（命题~\ref{prop:P1dead}：冠军几何上唯一可行的证书就是 $\rho_0I_3$）；
最后给出一个新的结构分解（定理~\ref{thm:section}），
把主线从「一个维数无界的帧流形上的上确界」改写成
「一个\textbf{二维单纯形}上的一族子问题」，并且指出冠军恰好落在其中的
\textbf{单轴}子族上。

顺带记录两条被定量否证的放缩族（命题~\ref{prop:ceilweyl}\ref{prop:ceilharm}），
它们分别给出 $\varphi=1.618034$ 与 $1.585642$ 的天花板，
都\textbf{劣于}报告已有的 $\Phi=\frac{3+\sqrt5}4=1.309017$，故只作为死路记录，不作为上界成果。

\subsection{凹性的解析证明与闭式梯度}\label{sec:r19:conc}

凹性在第~\ref{sec:r18}~节只有数值证据。它的解析证明其实只需要一行恒等式，
而这一行恒等式同时给出对偶族的显式形式，是本轮一切结论的起点。

\begin{lemma}[几何平均的仿射下确界表示]\label{lem:geo}
对 $a,b\ge0$，
\begin{equation}\label{eq:geoam}
\sqrt{ab}=\min_{t>0}\frac{ta+b/t}{2},
\end{equation}
极小值在 $t=\sqrt{b/a}$ 处取到（$a=0$ 或 $b=0$ 时下确界为 $0$，在 $t\to\infty$ 或 $t\to0$ 时逼近）。
故 $(a,b)\mapsto\sqrt{ab}$ 是 $\R_+^2$ 上一族仿射函数的逐点下确界，从而是\textbf{凹}的。
\end{lemma}

\begin{proof}
AM--GM 给 $\sqrt{ab}\le\frac{ta+b/t}2$，取 $t=\sqrt{b/a}$ 时等号成立。
一族仿射函数的下确界是凹函数，这是凸分析的标准事实。
\end{proof}

\begin{theorem}[$\rho(\Lam)$ 的凹性与几何平均表示]\label{thm:geo}
令 $D_i:=g_iu_iu_i^\top$、$L_i(\Lam):=\langle\Lam,D_i\rangle=g_iu_i^\top\Lam u_i=x_i(\Lam)^2$。
则
\begin{equation}\label{eq:geo}
\rho(\Lam)=\sum_{i,j}(Q_+)_{ij}\sqrt{L_i(\Lam)L_j(\Lam)}
=\inf_{T>0}\ \langle\Lam,S(T)\rangle ,
\end{equation}
其中 $T=(t_{ij})$ 取遍正对称矩阵，
\begin{equation}\label{eq:ST}
S(T):=\frac12\sum_{i,j}(Q_+)_{ij}\bigl(t_{ij}D_i+t_{ij}^{-1}D_j\bigr)\in\mathrm{Sym}_3 ,
\end{equation}
且当 $x(\Lam)>0$ 时下确界是\textbf{最小值}，在 $t^*_{ij}=x_j(\Lam)/x_i(\Lam)$ 处取到。
特别地 $\rho$ 在 $\{\Lam\succeq0\}$ 上\textbf{凹}且一次正齐次。
\end{theorem}

\begin{proof}
由 \eqref{eq:hlam}，$x_i(\Lam)=\sqrt{g_i}h_i(\Lam)=\sqrt{g_iu_i^\top\Lam u_i}=\sqrt{L_i(\Lam)}$，
代入 \eqref{eq:rholam}~的 $\rho(\Lam)=x(\Lam)^\top Q_+x(\Lam)$ 即得 \eqref{eq:geo}~的第一个等号。
对每一对 $(i,j)$ 用引理~\ref{lem:geo}（权 $(Q_+)_{ij}\ge0$）：
\[
\sum_{ij}(Q_+)_{ij}\sqrt{L_iL_j}
=\inf_{T>0}\frac12\sum_{ij}(Q_+)_{ij}\bigl(t_{ij}L_i+t_{ij}^{-1}L_j\bigr)
=\inf_{T>0}\Bigl\langle\Lam,\ \frac12\sum_{ij}(Q_+)_{ij}(t_{ij}D_i+t_{ij}^{-1}D_j)\Bigr\rangle ,
\]
这里用了 $L_i=\langle\Lam,D_i\rangle$ 的线性性。$L_i$ 是 $\Lam$ 的线性函数，
$\sqrt{\cdot\,\cdot}$ 凹，非负系数求和保凹，故 $\rho$ 凹；
一次正齐次性已见于定理~\ref{thm:lamparam}(i)。
\end{proof}

\begin{remark}[凹性为什么此前看起来「不好证」]\label{rmk:concwhy}
第~\ref{sec:r18}~节的注~\ref{rmk:concave}~指出凹性不能由 $\nabla^2\rho$ 的符号读出，
因为 \eqref{eq:gradlam}~只在 Perron 点成立。定理~\ref{thm:geo}~说明正确的看法是：
不要去算二阶导，而是先把 $\rho$ 写成\textbf{线性函数的几何平均的非负组合}。
这个表示的代价是引入了 $N(N+1)/2$ 个辅助参数 $t_{ij}$，
但回报是 $\rho$ 变成了一个\textbf{下确界型}（因此凹、因此可对偶）的泛函。
数值复核：冠军几何上 \eqref{eq:geo}~两种写法的最大差为 $6.7\times10^{-16}$，
$400$ 组随机 $(\Lam_1,\Lam_2,s)$ 的凹性余量最小值 $+1.2\times10^{-4}$，无一为负。
\end{remark}

\begin{corollary}[闭式梯度恒半正定]\label{cor:gradpsi}
在 $x(\Lam)>0$ 处 $\rho$ 可微，且
\begin{equation}\label{eq:gradpsi}
\nabla_\Lam\rho(\Lam)=S(T^*(\Lam))=\sum_i\frac{(Q_+x(\Lam))_i}{x(\Lam)_i}\,D_i
=\sum_i\frac{g_i[\langle u_i,m(\Lam)\rangle+d_i(\Lam)]}{h_i(\Lam)}u_iu_i^\top\ \succeq\ 0 .
\end{equation}
即梯度\textbf{恒落在锥} $\operatorname{cone}\{D_1,\dots,D_N\}$ \textbf{中}，
并且满足 Euler 恒等式 $\langle\Lam,\nabla_\Lam\rho(\Lam)\rangle=\rho(\Lam)$。
\end{corollary}

\begin{proof}
由 \eqref{eq:geo}~与 Danskin 定理（下确界在 $T^*$ 处唯一取到，因 AM--GM 的取等点唯一），
$\nabla_\Lam\rho=S(T^*(\Lam))$。把 $t^*_{ij}=x_j/x_i$ 代入 \eqref{eq:ST}：
\[
S(T^*)=\frac12\sum_i D_i\frac1{x_i}\sum_j(Q_+)_{ij}x_j+\frac12\sum_j D_j\frac1{x_j}\sum_i(Q_+)_{ij}x_i
=\sum_i\frac{(Q_+x)_i}{x_i}D_i ,
\]
第二步用 $Q_+$ 对称。$D_i\succeq0$ 与 $(Q_+x)_i/x_i\ge0$ 给出半正定。
第二个等号：$(Q_+)_{ij}=\sqrt{g_ig_j}(\langle u_i,u_j\rangle)_+$ 与 $x_j=\sqrt{g_j}h_j$ 给出
$(Q_+x)_i/\sqrt{g_i}=\sum_jg_jh_j(\langle u_i,u_j\rangle)_+
=\langle u_i,m\rangle+d_i$（把 $(\langle u_i,u_j\rangle)_+$ 拆成
$\langle u_i,u_j\rangle+(\langle u_i,u_j\rangle)_-$），再除以 $x_i=\sqrt{g_i}h_i$。
Euler 恒等式是一次正齐次性的标准推论。
\end{proof}

\begin{remark}[这条梯度公式的用处]\label{rmk:gradpsi}
其一，它把 \eqref{eq:gradlam}~从「临界点上的巧合」提升为「处处成立的闭式」，
因此\textbf{不需要}特征值或 Perron 向量就能求梯度，数值上可直接驱动内层凸规划
（本轮证书用它做 L-BFGS，$15$ 个起点全部收敛到 $\rho$，超出量 $\le1.3\times10^{-15}$）。
其二，$\nabla\rho(\Lam)\succeq0$ 说明 $\rho$ 关于 Loewner 序\textbf{单调}：
$\Lam_1\succeq\Lam_2\Rightarrow\rho(\Lam_1)\ge\rho(\Lam_2)$。
其三，$\nabla\rho$ 恒在 $\operatorname{cone}\{D_i\}$ 中，而 $\sum_iD_i=I_3$，
故 $\tr\nabla_\Lam\rho(\Lam)=\sum_ig_i(Q_+x(\Lam))_i/x(\Lam)_i$，
这是一个只涉及帧数据的显式标量；它与 $\rho(\Lam)$ 的关系由 Euler 恒等式给出。
\end{remark}

\subsection{精确对偶：$\Lam$-参数化没有松弛间隙，最优证书恒为各向同性}\label{sec:r19:dual}

\begin{theorem}[精确对偶]\label{thm:exactdual}
设 $x>0$（即 $Q_+$ 不可约；一般情形按锐图的分块取最大值化归）。令
$t^*_{ij}:=x_j/x_i$，其中 $x$ 是 $Q_+$ 的 Perron 向量。则
\begin{equation}\label{eq:stariso}
S(T^*)=\rho\,I_3 ,
\end{equation}
并且
\begin{equation}\label{eq:dualtight}
\max_{\substack{\Lam\succeq0\\ \tr\Lam=1}}\rho(\Lam)
\ =\ \inf_{T>0}\lambda_{\max}\!\bigl(S(T)\bigr)
\ =\ \lambda_{\max}(Q_+)\ =\ \rho .
\end{equation}
等式 \eqref{eq:stariso}~对\textbf{每一个}几何都成立，不需要任何临界点或极值假设。
\end{theorem}

\begin{proof}
\eqref{eq:stariso}：如推论~\ref{cor:gradpsi}~的证明，$S(T^*)=\sum_i\frac{(Q_+x)_i}{x_i}D_i$。
$x$ 是 $Q_+$ 的 Perron 向量，故 $Q_+x=\rho x$，于是 $\frac{(Q_+x)_i}{x_i}=\rho$，
$S(T^*)=\rho\sum_iD_i=\rho I_3$（帧条件 \eqref{eq:frame18}）。
\eqref{eq:dualtight}：三段夹逼。
(a) 弱对偶：对一切 $T$ 与一切 $\Lam\succeq0$，由 \eqref{eq:geo}~有
$\rho(\Lam)\le\langle\Lam,S(T)\rangle\le\lambda_{\max}(S(T))\tr\Lam$，
故 $\sup_{\tr\Lam=1}\rho(\Lam)\le\inf_T\lambda_{\max}(S(T))$。
(b) 由 \eqref{eq:stariso}，$\inf_T\lambda_{\max}(S(T))\le\lambda_{\max}(S(T^*))=\rho$。
(c) 由定理~\ref{thm:lamparam}(i)，$\sup_{\tr\Lam=1}\rho(\Lam)\le\lambda_{\max}(Q_+)=\rho$。
三者合起来迫使全部取等。
\end{proof}

\begin{remark}[这条定理说的是「坏消息」]\label{rmk:dualtight}
\eqref{eq:dualtight}~说明 $\Lam$-参数化\textbf{不产生任何松弛间隙}：
它的对偶最优值恰好是 $\lambda_{\max}(Q_+)$，也就是原问题本身。
这是好事也是坏事。好事是：定理~\ref{thm:lamparam}(iii)~的等价重构是\textbf{精确}的，
不存在「凸化之后丢掉信息」的隐患，因此 \eqref{eq:starmain}~确实值得作为主线。
坏事是：对偶证书路线\textbf{不能}给出比 $\lambda_{\max}(Q_+)$ 更好的东西——
下一小节把这句话变成对 P1 的否证。
\end{remark}

\begin{corollary}[内层规划是凸的，$\tilde\Lam$ 是全局极大点]\label{cor:argmax}
设 $x>0$。则对一切 $\Lam\succeq0$，
\begin{equation}\label{eq:tangent}
\rho(\Lam)\ \le\ \rho\cdot\tr\Lam ,
\end{equation}
从而 $\max_{\Lam\succeq0,\tr\Lam=1}\rho(\Lam)=\rho(\tilde\Lam)=\rho$。
于是定理~\ref{thm:lamparam}(iii)~的「$\ge$」一侧在 $x>0$ 时\textbf{无条件}成立
（第~\ref{sec:r18}~节里它被列为条件性结论）。
极大点集恰为仿射面
\begin{equation}\label{eq:face}
\mathcal F=\{\Lam\succeq0:\ \langle\Lam,D_i\rangle=x_i^2,\ i=1,\dots,N\},
\end{equation}
它退化为单点 $\{\tilde\Lam\}$ 当且仅当 $\{D_i\}_{i=1}^N$ 在 $\mathrm{Sym}_3$ 中线性无关
（等价地，那个 $N\times6$ 矩阵的秩为 $6$）。
\end{corollary}

\begin{proof}
$\rho$ 凹（定理~\ref{thm:geo}）且在 $\tilde\Lam$ 处可微、$\nabla\rho(\tilde\Lam)=\rho I_3$
（定理~\ref{thm:lamparam}(ii)），故一阶条件给
$\rho(\Lam)\le\rho(\tilde\Lam)+\langle\rho I_3,\Lam-\tilde\Lam\rangle=\rho+\rho(\tr\Lam-1)$。
对 $\tr\Lam\le1$ 得 $\rho(\Lam)\le\rho$；一般 $\Lam$ 用一次正齐次性化到 $\tr\Lam=1$，
即得 \eqref{eq:tangent}。
取等条件：\eqref{eq:geo}~中等号要求对每一对 $(Q_+)_{ij}>0$ 都有
$t^*_{ij}=\sqrt{L_j/L_i}$，即 $L_i/L_j=x_i^2/x_j^2$；锐图连通故存在 $c>0$ 使 $L_i=cx_i^2$。
而 $\tr\Lam=\sum_iL_i=c\sum_ix_i^2=c$，于是 $\tr\Lam=1$ 时 $c=1$，$L_i=x_i^2$，即 $\Lam\in\mathcal F$。
反过来 $\Lam\in\mathcal F$ 显然取等。$\mathcal F$ 是 $\mathrm{Sym}_3$ 中一个仿射子空间与半正定锥之交，
该仿射子空间由 $\{D_i\}$ 的正交补决定，故为单点当且仅当 $\{D_i\}$ 张成 $\mathrm{Sym}_3$。
\end{proof}

\begin{remark}[数值复核，以及一个意外发现]\label{rmk:argmax}
$\|\nabla\rho(\tilde\Lam)-\rho I_3\|_{\max}=1.6\times10^{-15}$；
$\rho(\tilde\Lam)-\rho=1.3\times10^{-15}$；
以 \eqref{eq:gradpsi}~为梯度的 L-BFGS 从 $15$ 个起点出发，最大超出量 $1.3\times10^{-15}$，
从未越过 $\rho$。切向界 \eqref{eq:tangent}~在 $960$ 个随机 $\Lam$（含 $\tr\Lam\ne1$）上最大违背 $0$。
\textbf{意外发现}：极大点\textbf{不总是唯一}的。$N=6$ 的一个上升极大元处
$\operatorname{rank}\{D_i\}=4<6$，此时 $\mathcal F$ 是一个二维仿射面，
L-BFGS 收敛到的极大点与 $\tilde\Lam$ 相距 $1.2\times10^{-1}$ 而\textbf{取值完全相同}（差 $1.1\times10^{-15}$）；
$N=7,9,10$ 处 $\operatorname{rank}\{D_i\}=6$，距离分别为 $0$、$1.2\times10^{-5}$、$2.1\times10^{-10}$。
这印证了 \eqref{eq:face}~的秩判据：\textbf{值}是刚性的，\textbf{极大点}只在 $\{D_i\}$ 线性相关时才不刚性。
\end{remark}

\subsection{各向异性对偶证书不存在：P1 被否证}\label{sec:r19:p1dead}

第八版队列的队首 P1 是：按注~\ref{rmk:certsearch}，构造形如
$S=\rho_0u_0u_0^\top+\kappa(I_3-u_0u_0^\top)$ 的\textbf{各向异性}证书。
本小节证明这件事在冠军几何上\textbf{不可能}，因此 P1 必须弹出。
证明只需要一条关于 $\tilde\Lam$ 的精确谱信息。

\begin{lemma}[冠军的 $\tilde\Lam$ 是单轴的，且正定]\label{lem:lampos}
冠军几何在绕枢纽轴的二面体群 $D_5$ 作用下不变，故 $\tilde\Lam$ 与之交换，从而
\begin{equation}\label{eq:uniax}
\tilde\Lam=\lambda_\parallel\,u_{\rm hub}u_{\rm hub}^\top+\lambda_\perp\bigl(I_3-u_{\rm hub}u_{\rm hub}^\top\bigr),
\qquad
\lambda_\parallel=h_{\rm hub}^2=\frac{\pi_{\rm hub}}{g_0^*},\quad
\lambda_\perp=\frac{1-h_{\rm hub}^2}{2},
\end{equation}
其中 $\pi_{\rm hub}=x_{\rm hub}^2=\dfrac{91-25\sqrt7}{126}=0.197271565265$，
$g_0^*=\dfrac{9-4\sqrt5+\sqrt{35}}{18}=0.331767104061$，故
$\lambda_\parallel=0.594608575866$、$\lambda_\perp=0.202695712067$。
特别地 $\tilde\Lam\succ0$，且这两个不等式是\textbf{精确}的：
$0<h_{\rm hub}^2<1$ 等价于 $28+28\sqrt5<25\sqrt7+7\sqrt{35}$，两边平方得 $6272<6090+350\sqrt{35}$，显然。
\end{lemma}

\begin{proof}
$\tilde\Lam$ 由命题~\ref{prop:uktkt}~的闭式 \eqref{eq:lamclosed}~从 $(g,U)$ 规范地定义，
故与 $D_5$ 作用交换。$D_5$ 在 $\R^3$ 上的表示分解为「轴上一维平凡」$\oplus$「轴外二维不可约」
（$n=5\ge3$ 时二维表示不可约），由 Schur 引理，交换代数中的元素在这两个不可约分量上是标量，
即 \eqref{eq:uniax}。$\lambda_\parallel=u_{\rm hub}^\top\tilde\Lam u_{\rm hub}=h_{\rm hub}^2$
（命题~\ref{prop:uktkt}(iv) 的 $u_i^\top\tilde\Lam u_i=h_i^2$）；
$\tr\tilde\Lam=1$（同处已证）给 $\lambda_\parallel+2\lambda_\perp=1$。
$h_{\rm hub}^2=x_{\rm hub}^2/g_{\rm hub}=\pi_{\rm hub}/g_0^*$ 是定义。
\end{proof}

\begin{proposition}[各向异性证书不存在]\label{prop:P1dead}
在冠军几何上，满足 \eqref{eq:certcond}~的 $S\in\mathrm{Sym}_3$ \textbf{只有一个}：$S=\rho_0I_3$。
因此推论~\ref{cor:dualcert}~的证书路线\textbf{等价于}原始不等式 \eqref{eq:starmain}，
不携带任何新信息；注~\ref{rmk:certsearch}~提出的各向异性候选族是空的。
\end{proposition}

\begin{proof}
设 $S$ 满足 \eqref{eq:certcond}。则
\[
\rho_0=\rho(\tilde\Lam)\ \le\ \langle\tilde\Lam,S\rangle\ \le\ \lambda_{\max}(S)\,\tr\tilde\Lam
\ =\ \lambda_{\max}(S)\ \le\ \rho_0 ,
\]
第一式是定理~\ref{thm:lamparam}(ii)~的 $\rho(\tilde\Lam)=\rho=\rho_0$，
第三式是引理~\ref{lem:lampos}~的 $\tr\tilde\Lam=1$。故三处全部取等。
由 $\langle\tilde\Lam,S\rangle=\lambda_{\max}(S)\tr\tilde\Lam$ 知 $\tilde\Lam$ 的像含于 $S$ 的
$\lambda_{\max}(S)=\rho_0$ 特征空间；而引理~\ref{lem:lampos}~给 $\tilde\Lam\succ0$，
其像为全空间，故该特征空间 $=\R^3$，$\lambda_{\max}$ 的重数为 $3$，$S=\rho_0I_3$。
\end{proof}

\begin{remark}[否证的定量内容]\label{rmk:p1dead}
这条否证有两层定量含义。
\textbf{第一层}：$S$ 的可行集是\textbf{单点}，不是「很小」；
任何试图调节 $\kappa$ 的做法都会立刻破坏 $\rho(\Lam)\le\langle\Lam,S\rangle$。
数值上，$3000$ 次随机各向异性尝试（$\lambda_{\max}(S)\le\rho_0$）中最好的一个仍然违反
证书条件达 $-4.2\times10^{-1}$。
\textbf{第二层}（更本质）：由定理~\ref{thm:exactdual}~的 \eqref{eq:stariso}，
\textbf{每个}几何的 AM--GM 取等参数都是 $t^*_{ij}=x_j/x_i$，都给出 $S=\rho I_3$。
所以对偶族 $\{S(T)\}$ 里「好的那一个」永远是各向同性的；
各向异性只能来自\emph{非}取等的 $T$，而那样 $\lambda_{\max}(S(T))>\rho$，不可能 $\le\rho_0$
（在冠军处 $\rho=\rho_0$）。
\textbf{结论}：凸对偶这条路在冠军几何上已经走到头，
主线必须回到不等式 \eqref{eq:starmain}~本身，或者改变分解方式。
\end{remark}

\subsection{两条被定量否证的放缩族}\label{sec:r19:ceil}

按协议，被排除的路线必须留下定量天花板。本轮排除两条自然但无效的族。

\begin{proposition}[Weyl$+$Frobenius 自举的天花板]\label{prop:ceilweyl}
记 $Q=VV^\top$、$Q_\pm=(\pm Q)_+$、$\alpha=\|V^\top x\|^2$、$\beta=x^\top Q_-x$，
则 $\rho=\alpha+\beta$，并且只用下面四条数据
\begin{equation}\label{eq:bootdata}
\tr Q_+=\textstyle\sum_ig_i=3,\qquad
\|Q_+\|_F^2+\|Q_-\|_F^2=\|Q\|_F^2=\tr(Q^2)=3,\qquad
\alpha\le\lambda_{\max}(Q)=1,\qquad \beta\le\lambda_{\max}(Q_-)\le\|Q_-\|_F ,
\end{equation}
最多只能推出
\begin{equation}\label{eq:boot}
\rho\ \le\ \frac{1+\sqrt5}{2}=\varphi=1.618033988750 ;
\end{equation}
若再加上 $N=6$ 时的两条迹精细化
$\|Q_+\|_F^2\ge\rho^2+\frac{(3-\rho)^2}{5}$（因 $\tr Q_+=3$）与
$\|Q_-\|_F^2\ge\frac65\lambda_{\max}(Q_-)^2$（因 $\tr Q_-=0$），则最多推出 $\rho\le\frac32$。
两者都\textbf{劣于}报告已有的 $\Phi=\frac{3+\sqrt5}4=1.309017$，故此族死。
\end{proposition}

\begin{proof}
$\rho=\alpha+\beta$ 是第~\ref{sec:r18}~节的分解。由 \eqref{eq:bootdata}，
$\rho\le1+\beta\le1+\|Q_-\|_F$ 且 $\rho^2\le\|Q_+\|_F^2=3-\|Q_-\|_F^2$，
消去 $\|Q_-\|_F$ 得 $(\rho-1)^2\le3-\rho^2$，即 $\rho^2-\rho-1\le0$，故 $\rho\le\varphi$。
$N=6$ 版本：$3\ge\rho^2+\frac{(3-\rho)^2}5+\frac65(\rho-1)^2$，
整理为 $2.4\rho^2-3.6\rho\le0$，即 $\rho\le\frac32$。
天花板是\textbf{可达}的（不是估计松紧的问题）：把 \eqref{eq:bootdata}~当约束做线性/二次规划，
$(\alpha,\beta)=(1,\varphi-1)$ 与 $\|Q_-\|_F^2=3-\rho^2$ 恰好相容，故 $\varphi$ 是该族的真上确界。
数值：$9$ 个阶数（$N=4,\dots,14$）的上升极大元全部满足 $(\rho-1)^2\le3-\rho^2$，
最大余量 $-1.28$；该不等式本身远未饱和，说明损失发生在 \eqref{eq:bootdata}~的第四条。
\end{proof}

\begin{proposition}[Legendre/优占族的天花板]\label{prop:ceilharm}
把核 $(t)_+$ 在 $[-1,1]$ 上按 Legendre 多项式展开，
\begin{equation}\label{eq:legc}
(t)_+=\sum_{\ell\ge0}c_\ell P_\ell(t),\qquad
c_\ell=\frac{2\ell+1}{2}\int_0^1tP_\ell(t)\dd t ,
\end{equation}
则 $c_0=\frac14$、$c_1=\frac12$、$c_2=\frac5{16}$、$c_3=0$、$c_4=-\frac3{32}$、$c_5=0$、
$c_6=\frac{13}{256}$、$c_8=-\frac{17}{512}$、$c_{10}=\frac{49}{2048}$、$c_{12}=-\frac{75}{4096}$，
且 \textbf{$c_\ell=0$ 对一切奇数 $\ell\ge3$}（因 $(t)_+=\frac t2+\frac{|t|}2$，而 $|t|$ 只有偶次项）。
令 $A_\ell:=\bigl(\sqrt{g_ig_j}\,P_\ell(\langle u_i,u_j\rangle)\bigr)_{ij}$，
则每个 $A_\ell\succeq0$（$P_\ell(\langle u,v\rangle)$ 是 $S^2$ 上的正定核，加法公式），
$P_+=\sum_\ell c_\ell A_\ell$，并且
\begin{equation}\label{eq:A01}
\lambda_{\max}(A_0)=\sum_ig_i=3,\qquad A_1=Q,\quad \lambda_{\max}(A_1)=1 .
\end{equation}
于是逐项放缩给出
\[
\rho\le\sum_{\ell:\,c_\ell>0}c_\ell\lambda_{\max}(A_\ell),
\]
而右端在冠军几何上截断到 $\ell=24$ 时等于 $1.585642$，比 $\rho_0$ 高 $0.304074$。
故这一族（含第~\ref{sec:nogo}~节的优占矩阵方案）不可能达到 $\rho_0$。
\end{proposition}

\begin{proof}
$c_\ell$ 由 \eqref{eq:legc}~直接积分（本轮用符号计算逐个求值，全部是有理数）。
奇数 $\ell\ge3$ 时 $tP_\ell(t)$ 是偶函数且 $\int_{-1}^1tP_\ell(t)\dd t=\int_0^1tP_\ell+\int_{-1}^0tP_\ell=0$
（用 $P_\ell$ 奇、$tP_\ell$ 偶），而 $\int_{-1}^1tP_\ell=0$ 对 $\ell\ne1$ 由正交性给出，
故 $\int_0^1tP_\ell=0$。$A_\ell\succeq0$ 与 $P_+=\sum_\ell c_\ell A_\ell$ 是逐条目的恒等式。
$A_0=\sqrt g\sqrt g{}^\top$ 给 $\lambda_{\max}=3$；$A_1=(\sqrt{g_ig_j}\langle u_i,u_j\rangle)=VV^\top=Q$。
\end{proof}

\begin{remark}[这条否证解释了一个数字]\label{rmk:harm54}
\eqref{eq:A01}~里 $\lambda_{\max}(A_0)=3$ 与 $\lambda_{\max}(A_1)=1$ 是\textbf{精确}的，
于是「只保留 $\ell\le1$」的逐项天花板是 $\frac14\cdot3+\frac12\cdot1=\frac54$，
恰好等于已证的 $\kappa_3(4)=\frac54$。这不是巧合式的装饰：
它说明 $\frac54$ 这一层信息完全来自核 $(t)_+$ 的\textbf{常数项与线性项}，
也就是「各向同性 $+$ 投影本身」这两件事；
$\rho_0-\frac54=0.031568$ 必须由 $\ell\ge2$ 的高次调和项提供。
但在冠军处这些高阶项的实际贡献是 $+0.051727$（其中 $\ell=6$ 一项就占 $+0.061846$，
$\ell=4,8,12$ 三项共 $-0.021199$），而且级数收敛很慢
（截断到 $\ell=24$ 时重构值 $1.279057$，误差 $-2.5\times10^{-3}$）。
\textbf{这就是逐项放缩失败的机制}：正负项交错、单项幅度大于总和，
任何「丢掉负项」的动作都会立刻越过 $\rho_0$。
\end{remark}

\subsection{Bloch 截面分解：主线的新形状}\label{sec:r19:sect}

到目前为止，主线一直是「在一个维数随 $N$ 增长的帧流形上求上确界」。
定理~\ref{thm:lamparam}~引入了 $5$ 个新参数（$\Lam\succeq0$、$\tr\Lam=1$，即一个三维 Bloch 体）。
本小节指出：因为帧集合在 $O(3)$ 下不变，\textbf{可以先对 $\Lam$ 的谱分类，再对帧取上确界}，
于是主线变成一个\textbf{二维单纯形}上的族。

\begin{theorem}[截面分解]\label{thm:section}
对 $\mu=(\mu_1,\mu_2,\mu_3)\in\Delta_2:=\{\mu_i\ge0,\ \sum_i\mu_i=1\}$（降序排列）定义
\begin{equation}\label{eq:section}
\kappa_3(\mu):=\sup_{N\ge3}\ \sup_{(g,U)\ \text{满足}\ \eqref{eq:frame18}}\ \rho_{g,U}(\operatorname{diag}\mu) .
\end{equation}
则
\begin{equation}\label{eq:secteq}
\kappa_3=\sup_{\mu\in\Delta_2}\kappa_3(\mu),
\end{equation}
且 $\kappa_3(R\Lam R^\top)=\kappa_3(\Lam)$ 对一切 $R\in O(3)$，
故 $\kappa_3(\mu)$ 只依赖 $\mu$ 作为\textbf{无序}三元组，即 $\kappa_3$ 是 $\Delta_2$ 上 $S_3$-不变的函数。
此外冠军满足
\begin{equation}\label{eq:mustar}
\mu(\tilde\Lam)=\Bigl(h_{\rm hub}^2,\ \tfrac{1-h_{\rm hub}^2}2,\ \tfrac{1-h_{\rm hub}^2}2\Bigr)
=(0.594608575866,\ 0.202695712067,\ 0.202695712067),
\end{equation}
是\textbf{单轴}的（后两个分量相等），且 $\kappa_3(\mu(\tilde\Lam))\ge\rho_0$。
\end{theorem}

\begin{proof}
乘积集上的上确界可换序：
$\sup_{(g,U)}\max_{\Lam}\rho_{g,U}(\Lam)=\sup_{(g,U),\Lam}\rho_{g,U}(\Lam)=\sup_{\Lam}\sup_{(g,U)}\rho_{g,U}(\Lam)$。
帧条件 \eqref{eq:frame18}~在 $u_i\mapsto Ru_i$ 下不变，且
$h_i(g,RU,\Lam)=\sqrt{(Ru_i)^\top\Lam(Ru_i)}=h_i(g,U,R^\top\Lam R)$、
$\langle Ru_i,Ru_j\rangle=\langle u_i,u_j\rangle$，故
$\rho_{g,RU}(\Lam)=\rho_{g,U}(R^\top\Lam R)$。
把 $R$ 换为 $R^\top$ 并对 $U$ 取上确界（$U\mapsto R^\top U$ 是帧集到自身的双射）即得不变性。
于是 $\kappa_3(\Lam)$ 是 $O(3)$-不变量，只依赖 $\operatorname{spec}(\Lam)$；
取 $\Lam=\operatorname{diag}\mu$ 得 \eqref{eq:secteq}。
\eqref{eq:mustar}~是引理~\ref{lem:lampos}；$\rho_{g^*,U^*}(\tilde\Lam)=\rho_0$
（定理~\ref{thm:lamparam}(ii)），故冠军是 $\mu(\tilde\Lam)$ 截面上的可行点。
\end{proof}

\begin{corollary}[两个特殊截面的显式形式]\label{cor:sectspecial}
\begin{enumerate}\itemsep3pt
\item[(i)] \textbf{中心截面} $\mu=(\frac13,\frac13,\frac13)$：$h_i\equiv\frac1{\sqrt3}$，故
\begin{equation}\label{eq:center}
\kappa_3\Bigl(\tfrac13,\tfrac13,\tfrac13\Bigr)=\frac13\sup_{(g,U)}\sum_{i,j}g_ig_j\bigl(\langle u_i,u_j\rangle\bigr)_+ ,
\end{equation}
即「各向同性测度 $\mu=\sum_ig_i\delta_{u_i}$ 的正部二次能量」除以 $3$。
此时 $\Lam$ 的自由度\textbf{完全消失}，问题变成一个纯测度论的极值问题。
\item[(ii)] \textbf{秩一截面} $\mu=(1,0,0)$，$\Lam=ww^\top$：$h_i=|\langle u_i,w\rangle|$，
令 $a_i:=\sqrt{g_i}\langle u_i,w\rangle$，则 $\sum_ia_i^2=\|w\|^2=1$ 且 $x_i(\Lam)=|a_i|$，故
$\rho_{g,U}(\Lam)=\sum_{ij}(Q_+)_{ij}|a_i||a_j|$。
\end{enumerate}
\end{corollary}

\begin{proposition}[截面地图，数值]\label{prop:sectmap}
对 $8$ 个截面 $\mu$ 与 $N\in\{6,8,10\}$，用 \eqref{eq:gradpsi}~之外的直接梯度上升
（对 $V\in\mathrm{St}(N,3)$，解析梯度）求 $\kappa_3(\mu)$ 的数值下界，得
\begin{center}\small
\begin{tabular}{lrrr}
\toprule
截面 $\mu$ & $N=6$ & $N=8$ & $N=10$ \\
\midrule
中心 $(\frac13,\frac13,\frac13)$ & $1.233669$ & $1.256733$ & $1.248877$ \\
单轴 $a=0.50$ & $1.249734$ & $1.249025$ & $1.245552$ \\
单轴 $a=0.5946=h_{\rm hub}^2$ & $1.250166$ & $1.270937$ & $1.273641$ \\
单轴 $a=0.70$ & $1.247226$ & $1.240466$ & $1.278518$ \\
单轴 $a=0.85$ & $1.209759$ & $1.222580$ & $1.261757$ \\
秩二 $(\frac12,\frac12,0)$ & $1.228572$ & $1.181850$ & $1.182550$ \\
秩一 $(1,0,0)$ & $1.150694$ & $1.175367$ & $1.166373$ \\
非对称 $(0.60,0.30,0.10)$ & $1.231907$ & $1.261213$ & $1.260940$ \\
\bottomrule
\end{tabular}
\end{center}
其中「单轴 $a$」指 $\mu=(a,\frac{1-a}2,\frac{1-a}2)$。
把冠军本身作为 $a=h_{\rm hub}^2$ 截面的可行点加入后，该截面的值为 $1.281568196872167=\rho_0$。
\textbf{观察}：（1）没有任何截面的数值上确界超过 $\rho_0$；
（2）最大值出现在单轴族 $a\approx0.6$--$0.7$，冠军恰在其中；
（3）秩一与秩二截面明显更低（最大 $1.228572$，余量 $0.052996$）；
（4）非对称截面 $(0.60,0.30,0.10)$ 也不超过 $1.2613$。
这些都是\textbf{数值}结论，不作为成果。
\end{proposition}

\begin{remark}[这条分解买到了什么]\label{rmk:sectread}
\textbf{买到三件事}。
其一，$\Lam$ 的 $8$ 个自由度被压缩成 $\Delta_2$ 的 $2$ 个：
旋转自由度是\textbf{冗余}的（定理~\ref{thm:section}~的不变性），
这是一个此前没有明确使用过的事实。
其二，$\Delta_2$ 的\textbf{边界}（秩 $\le2$ 的截面）与\textbf{中心}截面各自成为独立的、
形式上更简单的子问题；边界截面数值上余量很大（$\ge0.053$），
中心截面按 \eqref{eq:center}~根本不含 $\Lam$，因此最有希望先被解析解决。
其三，冠军落在\textbf{单轴}子族上，这不是偶然：
引理~\ref{lem:lampos}~说明单轴性是 $D_5$ 对称性的后果，
而冠军的极值性本身就来自 $D_5$ 对称。
\textbf{代价}是：$\kappa_3(\cdot)$ 是「一族凹函数的上确界」，本身既不凹也不凸，
所以不能用凹性把它压到某个特殊截面上；
排除非对称极大元（$S_3$-轨道大小为 $6$ 的那些）需要新的论证。
\end{remark}

\begin{corollary}[归约的精确形式]\label{cor:sectreduce}
主猜想 $\kappa_3\le\rho_0$ 等价于：对每个 $\mu\in\Delta_2$，$\kappa_3(\mu)\le\rho_0$。
由 $S_3$-对称性，只需对下列三类 $\mu$ 证明：
\textbf{(A)} 中心 $\mu=(\frac13,\frac13,\frac13)$；
\textbf{(B)} 单轴 $\mu=(a,\frac{1-a}2,\frac{1-a}2)$，$a\in[\frac13,1]$；
\textbf{(C)} 真正的非对称极大元（若存在）。
若能证明 $\kappa_3(\cdot)$ 在 $\Delta_2$ 上的极大元必落在 (A)$\cup$(B) 中，
则主猜想化归为\textbf{一个单参数族}。
\end{corollary}

\subsection{本轮状态、诚实标签与方法优先队列（第九版）}\label{sec:r19:status}

\textbf{已证（解析，含完整证明）}：
引理~\ref{lem:geo}（几何平均的仿射下确界表示）；
定理~\ref{thm:geo}（$\rho(\Lam)$ 的凹性与 $\inf_T\langle\Lam,S(T)\rangle$ 表示）；
推论~\ref{cor:gradpsi}（闭式梯度、恒半正定、Loewner 单调、Euler 恒等式）；
定理~\ref{thm:exactdual}（$S(T^*)=\rho I_3$ 对每个几何成立；对偶精确紧）；
推论~\ref{cor:argmax}（切向界 \eqref{eq:tangent}、$\tilde\Lam$ 是全局极大点、
极大点面 \eqref{eq:face}~及其秩判据；定理~\ref{thm:lamparam}(iii)~的「$\ge$」一侧在 $x>0$ 时无条件）；
引理~\ref{lem:lampos}（冠军 $\tilde\Lam$ 单轴且正定，$0<h_{\rm hub}^2<1$ 的精确证明）；
命题~\ref{prop:P1dead}（各向异性对偶证书不存在）；
命题~\ref{prop:ceilweyl}（Weyl$+$Frobenius 自举的天花板 $\varphi$ / $\frac32$）；
命题~\ref{prop:ceilharm}（Legendre 系数精确值、$A_\ell\succeq0$、$\lambda_{\max}(A_0)=3$、$\lambda_{\max}(A_1)=1$、
逐项天花板 $1.585642$）；
定理~\ref{thm:section}（截面分解）；推论~\ref{cor:sectspecial}（中心与秩一截面的显式形式）；
推论~\ref{cor:sectreduce}（三类归约）。

\textbf{条件性}：定理~\ref{thm:exactdual}~与推论~\ref{cor:argmax}~假设 $x>0$
（即 $Q_+$ 不可约 / 锐图连通）；一般情形的化归（按分块取最大值）在本报告中按
\textbf{条件性}处理，与第~\ref{sec:r18}~节 $g_i=0$、重根的退化情形是同一类技术债务。
推论~\ref{cor:sectreduce}~的 (C) 类（排除非对称极大元）尚未解决。

\textbf{数值（不作成果）}：
命题~\ref{prop:sectmap}（$8$ 个截面 $\times3$ 个阶数的地图）；
定理~\ref{thm:exactdual}~的 $13$ 个几何复核（$|S(T^*)-\rho I_3|_{\max}\le9.1\times10^{-15}$，
弱对偶最大违背 $0$）；推论~\ref{cor:argmax}~的 L-BFGS 复核（超出 $\le1.3\times10^{-15}$）；
命题~\ref{prop:ceilweyl}~的 $9$ 个阶数复核；命题~\ref{prop:ceilharm}~的 $\ell\le24$ 截断复核。
证书脚本：\texttt{证书\_凹性对偶与截面分解.py}（$23$ 项全绿）。

\textbf{否证}：命题~\ref{prop:P1dead}（各向异性对偶证书；可行集是单点 $\{\rho_0I_3\}$，
随机搜索最好者仍违反 $-4.2\times10^{-1}$）；
命题~\ref{prop:ceilweyl}（天花板 $\varphi=1.618034$，比 $\rho_0$ 高 $0.336466$）；
命题~\ref{prop:ceilharm}（天花板 $1.585642$，比 $\rho_0$ 高 $0.304074$）。
这三条与第~\ref{sec:r18}~节的命题~\ref{prop:refA}\ref{prop:refB}~一起，
把「凸化 / 对偶 / 逐项放缩」这一整片方法\textbf{全部}排除在主线路径之外。

\textbf{方法优先队列（第九版）}
\begin{enumerate}\itemsep3pt
\item \textbf{P1（队首，全局）——单轴截面归约}。目标：证明推论~\ref{cor:sectreduce}~的 (C) 类不存在，
即 $\kappa_3(\cdot)$ 在 $\Delta_2$ 上的极大元必为中心或单轴点。
入口（按优先级）：
(a) \textbf{对称临界点论证}：$\kappa_3$ 是 $S_3$-不变的；若它在 $\mu$ 处可微，
则 $\nabla\kappa_3(\mu)$ 必与稳定子群的不动点空间正交，
故稳定子群非平凡的点（中心、单轴）自动是临界点；需要的是\emph{排除}稳定子群平凡的极大元。
(b) \textbf{凹性的替代物}：$\kappa_3$ 是凹函数的上确界，本身不凹；
但对\textbf{固定}帧 $\rho_{g,U}(\cdot)$ 凹，故
$\kappa_3(\mu)\ge\int\rho_{g,U}(R\operatorname{diag}(\mu)R^\top)\dd R$ 对每个帧成立，
右端是 $\mu$ 的\emph{对称化凹}函数。若能证明对称化不改变上确界的位置，则 (C) 排除。
这一步目前\textbf{没有证明}，是第一障碍。
(c) 退一步：不求排除 (C)，而是直接对\textbf{单轴族}做解析（见 P2），
把 (C) 的处理留到最后。
\item \textbf{P2（全局）——中心截面的精确值}。由 \eqref{eq:center}，这是
$\frac13\sup\sum_{ij}g_ig_j(\langle u_i,u_j\rangle)_+$，
即「$S^2$ 上各向同性测度的正部能量」。它\textbf{不含} $\Lam$，是本轮之后形式最干净的子问题。
数值上 $\kappa_3(\frac13,\frac13,\frac13)\approx1.256733$（$N=8$），
小于 $\rho_0$ 达 $0.024836$。已知工具只能给到 $\Phi=1.309017$
（用 $(\langle u,v\rangle)_+\le\frac{|\langle u,v\rangle|+\langle u,v\rangle}2$ 与 $\lambda_{\max}(|Q|)\le\varphi$）。
需要的是一条\textbf{正部特有的}不等式：正部只在锐角对上计数，
而各向同性条件 $\sum_ig_iu_iu_i^\top=I_3$ 限制了锐角对的总权重。
\item \textbf{P3（全局）——秩一/秩二截面}。数值余量最大（$\ge0.053$），
推论~\ref{cor:sectspecial}(ii)~已把它写成 $x_i=|a_i|$、$\sum a_i^2=1$ 的形式，
应当是最容易先关掉的一类。
\item \textbf{P4（全局）——$\mathcal E_-$ 的下界}（第八版的 P3，未动）：
用引理~\ref{lem:Em2}~的显式 $\zeta_{ij}$ 在 \eqref{eq:star}~中给 $\frac12\mathcal E_-\ge\Phi-\tau_0$。
\item \textbf{P5（局部）——第~\ref{sec:r1517}~节的分支 I / 分支 II}（第八版的 P4）。
定理~\ref{thm:transitive}~给出分支 I 的新入口。
\item \textbf{P6（$N=6$ 收尾）}：第~\ref{sec:n6global}~节余下的 $20$ 个符号类。
\item \textbf{P7（副方向，下界）}：$\Lam_3(G)\ge\rho_0n-3$ 的可达到性，只差定理~\ref{thm:Fn}~的无穷族渐近。
\end{enumerate}

\textbf{已弹出}：第八版的 P1（构造对偶证书 $S$）——命题~\ref{prop:P1dead}~证明其可行集是单点；
第八版的 P2（解析证明凹性）——定理~\ref{thm:geo}~已完成。

\textbf{本轮把主线的形状改成了}：
\[
\kappa_3=\sup_{\mu\in\Delta_2}\kappa_3(\mu),\qquad
\kappa_3(\mu)=\sup_{N}\ \sup_{(g,U)}\ \rho_{g,U}(\operatorname{diag}\mu),
\]
其中 $\rho_{g,U}$ 对每个帧都是\textbf{凹}的、有\textbf{闭式梯度}的、
其内层对偶\textbf{精确紧}的泛函。$\Lam$ 的 $8$ 个自由度里 $6$ 个是冗余的旋转，
只剩 $\Delta_2$ 的 $2$ 个。冠军坐在单轴子族上，
中心截面则退化成一个不含 $\Lam$ 的经典测度极值问题 \eqref{eq:center}。
凸对偶与逐项放缩这两片方法已经被完整排除（命题~\ref{prop:P1dead}%
\ref{prop:ceilweyl}\ref{prop:ceilharm}），
下一步必须用\textbf{各向同性条件与正部之间的直接相互作用}，而不是再找一层松弛。


